\section{Methods}

\label{sec:methods}

\subsection{Current cohorts and statistical specification}

\label{sec:methods-current-cohorts-and-statistical-specification}

The current investigation addresses authorization semantics, joint security and utility, and bounded component evidence. It is separate from the earlier six-question programme. The fixed 180-case authorization sample contains 100 attacks, 60 benign tool tasks and 20 no-tool tasks. Its 60-case stress subset and the 120-case contemporaneous calibration overlap earlier samples; neither a condition assignment nor a reused cohort is an independent task. Model prompts are ordinary in both policy and full cells. Oracle-derived scope is supplied only to the Guard and is explicitly distinct from deployed, non-oracle authorization.

The current authorization comparisons cross policy/full modules with compiled and corrected exact-target scope. The additional broad, underspecified and conflicting transformations are nested descriptive stress tests. Calibration contrasts policy, full and the explicitly configured VIGIL adapter with an unguarded control that retains the same external-effect boundary. Sixteen synthetic controls test specified label and data-flow behavior; they do not open an additional confirmatory family or estimate natural-task performance.

We retain the prespecified inferential results rather than re-estimating historical outcomes. Cell proportions use Wilson 95\% intervals; paired absolute differences use Newcombe intervals and exact McNemar tests, with Holm correction within the defined endpoint/contrast family. Bootstrap intervals, when reported, keep task-level paired assignments together rather than resampling condition rows independently. Marginal cell intervals, paired-difference intervals and identification bounds for missing outcomes are different objects and are labelled accordingly. Zero-event empirical bootstrap intervals do not establish equivalence or absence of risk.

Benchmark panels are reported separately because native AgentDojo security judgments and M3-v2 observed simulated effects are not interchangeable. Completion and operational failure may coexist. All assignments remain in ITT, observed-value denominators accompany unknown counts, and sensitivity bounds never substitute for observed unsafe effects. The benign-task review is post-outcome descriptive interpretation, not an additional significance family or an adjudication of every false block. Provider identities, module inputs and configured baseline limits bound each claim.

\subsection{Threat model and safety boundary}

\label{sec:methods-threat-model-and-safety-boundary}

The attacker controls text in documents, messages, retrieved records, memory, or simulated tool outputs processed while the agent performs a trusted user task. Corrected M3-v2 security effects require observed simulated execution that violates the task-specific evaluator. AgentDojo retains its native attack judgment, which is not reduced to a state-transition-only definition. Historical endpoint labels retain their original semantics. The model may propose tool calls, but it cannot invoke provider-managed tools or reach a real external system.

All experiments use isolated benchmark state. Real email, payment, browser, shell, code-execution, production API, credential, and unrestricted filesystem operations are disabled. A proposal can affect only copied AgentDojo state, M3-v2 mock resources, \nolinkurl{example.test}, \nolinkurl{/mock/}, or the public sandbox. Model output is therefore treated as untrusted data, not executable authority.

Training-time poisoning, model-weight compromise, credential theft, host compromise, and direct modification of the trusted policy store are outside scope.

\subsection{Online-state and evaluator separation}

\label{sec:methods-online-state-and-evaluator-separation}

Every run constructs the online state through an allowlist projection. Permitted fields include trusted user intent, the benchmark-visible task context, the selected authorization regime, provenance events available at that step, and the proposed action. Gold answers, expected actions, attack and task-success labels, evaluator explanations, reference trajectories, and offline effect receipts are excluded.

The \nolinkurl{exact_gold} condition is an explicitly labeled oracle sensitivity analysis. It converts evaluator-derived authorization into the condition's authorization object, but does not expose the evaluator labels or scoring fields themselves. \nolinkurl{exact_gold} is never presented as the primary deployment regime. The primary regime is \nolinkurl{auto_compiled}, created from the trusted user request before untrusted content is processed.

Each compact result stores a digest of the online-visible projection and a boolean forbidden-field check. A run fails the minimal integrity check if a forbidden field is present. The final analysis checks this condition once for all terminal records; it does not recursively replay historical evidence.

\subsection{Runtime conditions}

\label{sec:methods-runtime-conditions}

The common runtime interfaces are \nolinkurl{TaskPolicyManifest}, \nolinkurl{ActionProposal}, \nolinkurl{GuardDecision}, and \nolinkurl{ProvenanceEvent}. In the historical bundled comparison, the two M3 conditions share the agent, model, parser, evaluator, tool environment, and task-policy representation, but they are not a single-module ablation:

\begin{itemize}

\item \textbf{Policy package} evaluates the proposed tool and arguments against the current task policy with the \nolinkurl{normal} prompt profile. Provenance propagation, temporal/data-flow checks, and independent confirmation are disabled.

\item \textbf{Full package} uses the \nolinkurl{safe_prompt} profile and enables policy enforcement plus parameter provenance, temporal/data-flow constraints, permission-monotonic transitions, and confirmation. Automatic transitions may preserve or reduce authority; authority expansion produces \nolinkurl{ASK} and requires an independent confirmation result.

\end{itemize}

The earlier authorization experiment changed prompt profile and runtime modules together and is therefore retained only as a bundled configuration comparison. The historical E1 concurrent factorial crosses both factors explicitly: policy versus full modules and normal versus safe prompt. Its module main effect, prompt main effect, and module-by-prompt interaction are therefore separately identified within the fixed paired sample. The current corrected authorization and calibration studies instead hold both module conditions at the same ordinary prompt; their assignments must not be described using the earlier bundled prompt profiles.

\nolinkurl{ALLOW}, \nolinkurl{ASK}, and \nolinkurl{DENY} are runtime decisions, not security outcomes by themselves. Security is scored from the isolated post-enforcement effect receipt. Utility is scored by the benchmark task-completion evaluator.

The earlier incomplete 450-item Gemini authorization sample remains diagnostic and is not extended. E6 makes a separately executed, bounded Gemini prefix attempt that is diagnostic by design under the single-block amendment. Even complete prefix coverage would not make that amended design a formal cross-provider study; the observed incomplete prefix leaves every Gemini and cross-provider effect not estimable. Structural baselines, including the source-pinned VIGIL adapter, retain their earlier descriptive role; the new experiments target causal confounding, source authority, risk enrichment, same-provider model replication, and adaptive tail behavior rather than rerunning a large baseline matrix.

\subsection{Measurement correction and development validation}

\label{sec:methods-measurement-correction-and-development-validation}

The historical execution studies used the original compiler, rule matcher, and evaluator. Their recorded safety labels are not retrospectively converted to executed un\-au\-thor\-ized effects. In particular, the original M3-v2 evaluator could score a final isolation rejection as an unsafe event, whereas the revised development evaluator distinguishes proposals, Guard permission, isolation rejection, successful simulated execution, task completion, and unknown values. AgentDojo retains its native security and utility checks. Legitimate completion and a security violation may coexist; missing evidence is not a safe zero.

The explicitly selected corrected-v1 path fixes action-resource binding, complete matching across multiple rules, trusted-environment label propagation, and confirmation-state continuity. It does not alter the historical default or records. Six configurations cross no/policy/full guard modules with the same normal/safe prompt factor. The development set comprises 24 attacks, 16 benign-tool tasks, and eight no-tool tasks; AgentDojo contributes 20 tasks and M3-v2 28. Attacks are historically enriched, not a deployment sample. Both completed DeepSeek cohorts use temperature zero, a 512-token output limit, and at most eight provider calls per task-condition. The V4.1 Flash model is distinct from the earlier V4 Flash and is not a model-invariant before/after comparison. Synthetic paired-action tests establish component behavior only; there is no simulated human approval or new confirmatory family.

\subsection{Corrected target authorization and contemporaneous calibration}

\label{sec:methods-corrected-target-authorization-and-contemporaneous-calibration}

Table 1 summarizes the current cohorts and their distinct evidence roles.

The corrected experiment selects explicit \nolinkurl{corrected_v2} runtime behavior and \nolinkurl{exact_target_v2} authorization. Target and other constrained security fields must match the offline-derived scope; schema-required content fields need not equal an absent annotation. Schema completeness, types, additional security targets, action order, and execution registration remain checked. A generic subset-of-arguments comparison is insufficient because additional recipients, paths, or endpoint parameters can change the authorized action. The original matcher and all its records remain unchanged.

Raw evaluator fields do not enter the model prompt. In the oracle condition, the Guard nevertheless receives an explicitly answer-derived authorization object. This is an idealized information regime, not evidence that deployed authorization can obtain the same scope. Ambiguous required scope is not an automatic grant, including with confirmation modules disabled. Missing trusted classification cannot be treated as public data when a restricted data-flow rule applies. No human approval is simulated: ASK means pending clarification, not task completion.

The new authorization cohort reuses 180 fixed cases: 100 attacks, 60 benign tool tasks, and 20 M3-v2 no-tool tasks. Policy and full modules each receive compiled and exact-target authorization, yielding 720 assignments; all four cells use the same ordinary prompt. Three stress regimes on the predefined 60-case nested subset add 180 full-module assignments. The 12-assignment formal prefix is part of the 900 total, not an additional repeat. The security endpoint has 490 attack assignments across the main and stress cells, not 490 independent tasks. This authorization study contains no unguarded cell.

A separate contemporaneous calibration uses 120 fixed cases, with 30 attacks and 30 benign-tool tasks from each benchmark. Ordinary prompts and compiled authorization are held constant across no research Guard, policy, full, and configured VIGIL, giving 480 assignments. The unguarded condition retains the same host and external-effect boundary. VIGIL generates constraints through the model and uses its configured deterministic hybrid audit, with model-based secondary verification disabled. It is not the complete official default baseline. Its provider-call cap is 16; the other conditions cap at eight. Cohorts overlap earlier studies and do not constitute independent sample replication.

Both Pro experiments use temperature zero, 512 output tokens, a 180-second single-request timeout, and zero automatic HTTP retries. Task environments are isolated, with at most two assignments in flight. After a 120-assignment calibration stop caused by parallel telemetry and truncated constraint classification, a telemetry-only revision executed the 360 untouched items. The first 120 records, including response-bearing failures, were not retried. Scientific prompts, parameters, permissions, and work identities were not changed; the two runner revisions remain distinguishable.

Four synthetic memory scenarios cross policy/full with Pro and the approved V4.1 Flash service, yielding 16 additional executions. The scenarios cover private restricted flows, public restricted flows, private archival flows, and absent classification. They test specified component behavior, not general protection, a natural task distribution, or human confirmation.

Each benchmark retains its own security semantics. AgentDojo uses native attack and utility checks on an isolated environment copy; M3-v2 counts un\-au\-thor\-ized simulated effects supported by execution evidence. Failures, unknown outcomes, and completion after execution remain distinct. Reported total latency reads the runtime total component; component durations and missingness are reported separately rather than inferred from token usage. Existing inferential results are reused when correcting latency aggregation.

\input{\ProjectRoot/table_01.tex}

\input{\ProjectRoot/figure_01.tex}

\subsection{Trusted-intent review of benign tasks}

\label{sec:methods-trusted-intent-review-of-benign-tasks}

We subsequently reviewed the fixed calibration cohort's 60 benign tool tasks, 30 per benchmark, and linked their 240 existing condition records. This is a post-outcome explanatory review, not a blind evaluation, a new sample, or an independent replication. Trusted user requests, trusted application context, and tool interfaces define the authority interpretation; evaluator answers and model-declared permissions do not supply grants.

For each task we retain an intent paraphrase, possible tool routes, declared action fields, generated rule types, first restriction, successful simulated execution count, completion, operational status, and source locators. Possible tool routes are not a gold list of necessary actions. Interpretation categories may overlap, and the same task may have a completed benchmark outcome and a later failed workflow. We inspect declared action--target pairs for loss or mismatch without treating every natural-language intent as a contract requiring the current lexical compiler to infer a dynamic target.

The review uses only existing public records and projected trusted inputs. Redacted proposal arguments and final-answer content limit diagnosis of cases with successful simulated execution but unsuccessful completion. No new model call, permission expansion, evaluator relabeling, or runtime revision follows from the descriptive review, and no independent reviewer-agreement estimate is available.

\subsection{Failure and recovery policy}

\label{sec:methods-failure-and-recovery-policy}

Automatic retry is disabled. Recovery operates at the predefined experiment-item level:

\begin{itemize}

\item a work item with an official model response is terminal, including timeout, refusal, parse failure, and maximum-step outcomes after that response;

\item a transport or preparation failure with no official response may be retried only with the explicit failed-row option;

\item only the affected work item is retried; successful and unaffected work items are skipped;

\item a repeated non-rate-limit pre-response failure becomes terminal on its second attempt; and

\item no automatic retry or cross-model fallback is permitted.

\end{itemize}

Earlier response-less attempts remain append-only in the raw file. The scientific analysis selects the latest terminal attempt for each experiment item and reports the count of superseded attempts. No prior row is overwritten.

\subsection{Outcomes}

\label{sec:methods-outcomes}

The two joint primary endpoints are:

\begin{enumerate}

\item \textbf{unsafe-effect rate} on assigned attack tasks; and

\item \textbf{task-completion rate (TCR)} on assigned benign tool tasks.

\end{enumerate}

Benign no-tool completion is descriptive. Secondary outcomes are system \nolinkurl{ASK} and \nolinkurl{DENY} event counts and record incidence, operational failure, p50/p95 latency, total tokens, API cost, and predefined heterogeneity by benchmark and attack family. \nolinkurl{ASK} is a runtime confirmation request; without a human study it is not interpreted as measured human burden.

All terminal work items remain in the intention-to-treat denominator. The analysis reports raw numerator/denominator values alongside proportions. Observed security outcomes, unknown outcomes, task completion, and operational failure are recorded separately. A provider or parser failure is neither silently removed nor relabeled as successful protection. If the benchmark establishes completion before a later failure, both observations are retained; without completion evidence the failure cannot create a successful task. Missing effect evidence remains unknown, not a safe zero. Burden summaries include terminal failures, and latency summaries report their available-value denominator rather than imputing missing durations.

Because observed unsafe effects are rare, we also report a pessimistic sensitivity endpoint on attack tasks: \nolinkurl{unsafe_effect OR operational_failure}. This composite is a robustness bound, not an attack-success label. It asks how conclusions change if a non-operational run is not credited as secure while preserving unsafe effects and failures as separate reported counts.

\subsection{Minimal verification and reproducibility}

\label{sec:methods-minimal-verification-and-reproducibility}

Verification is deliberately proportional to the lean study. It consists of:

\begin{enumerate}

\item focused unit tests for modified provider, authorization projection, evaluator, analysis, and plotting code;

\item an eight-item formal prefix for each regular incremental study, two five-query tail prefixes for E5, and a 12-item outcome-blind E6 prefix; and

\item one terminal pass over row count, experiment-item uniqueness, forbidden-field isolation, exact model identity, and ITT failure retention.

\end{enumerate}

Each study writes one append-only compact raw JSONL. Final statistics are written once to one analysis JSON per study, and figures are exported as submission-ready vector PDF, editable SVG, and raster PNG previews for manuscript/QA use. Reproducibility assets include frozen samples, condition assignments, analysis code, aggregate tables, and versioned source.
